\section{The scalar determinant of the acyclotomic twist}
\label{sec:twist-scalar-determinant}

Lochak--Schneps express an acyclotomic GT element through four
block-substitution actions on a free profinite group of rank three
\cite[Theorem~2.1, equation~(8), pp.~121--122]{LS2015Twists}.
The definitions of the four actions are equations~(9)--(12),
and their peripheral product preservation is Lemma~3.1
\cite[pp.~122--124]{LS2015Twists}. We compute what their scalar
Fox determinants retain.

There is a printed error in equation~(10): the first conjugator in
the displayed image of \(z\) has first argument \(y\).
The block-substitution definition gives instead
\[
 F^{2,3}(z)=f(z,(yz)^{-1})\,z\,f((yz)^{-1},z).
\]
We use this block-substitution action. It is also the one compatible
with the two-generator action and the product assertion of
Lemma~3.1 in the cited source.

Use \([x,y]=xyx^{-1}y^{-1}\), and write
\[
 [f]=R_f(X,Y)[x,y],\qquad
 B_f(X,Y)=1+(X-1)(Y-1)R_f(X,Y).
\]
All Fox derivatives are abelianized. They can be computed in the
completed coefficient-two group algebra and then specialized to
finite abelian-cover quotients. For an IA automorphism \(\alpha\)
of a free profinite group, let \(J_\alpha\) be the matrix of
abelianized Fox derivatives, with the image of the \(i\)-th
generator in row \(i\). Its determinant is multiplicative on IA
automorphisms, since their action on the coefficient algebra is
trivial.

\begin{proposition}\label{prop:twist-determinant-beta}
For the four block-substitution actions of an acyclotomic element
\((1,f)\), the Fox determinants are respectively
\[
 B_f(X,Y),\qquad B_f(Y,Z),\qquad
 B_f(XY,Z),\qquad B_f(X,YZ)
\]
for \(F^{1,2},F^{2,3},F^{12,3},F^{1,23}\).
The action on the rank-three fiber of the pure-braid proword
\(f(x_{12},x_{23})\) has determinant one. Consequently, at
the determinant level the twist relation gives exactly
\begin{equation}
 B_f(X,Y)B_f(XY,Z)=B_f(X,YZ)B_f(Y,Z).
 \label{eq:twist-determinant-beta}
\end{equation}
This determinant operation supplies no scalar equation beyond the
beta two-cocycle identity.
\end{proposition}

\begin{proof}
First consider \(\phi_f(x)=x\), \(\phi_f(y)=f^{-1}yf\).
The abelianized derivative vector of \(f\) is
\[
 (\partial_x f,\partial_y f)
          =R_f(X,Y)(1-Y,X-1).
\]
Thus
\[
 J_{\phi_f}=
 \begin{pmatrix}
 1&0\\
 -(Y-1)^2R_f&1+(X-1)(Y-1)R_f
 \end{pmatrix},
 \qquad \det J_{\phi_f}=B_f(X,Y).
\]
If \(w\) belongs to the commutator subgroup of a free group of
rank \(d\), and \(q_j=\partial_j w\), then
\[
 J_{\operatorname{Inn}(w)}
   =I+\bigl(1-X_i\bigr)_i(q_j)_j.
\]
The determinant lemma and the Fox fundamental identity give
\[
 \det J_{\operatorname{Inn}(w)}
 =1+\sum_j(1-X_j)q_j=1,
\]
because the abelianization of \(w\) is one. An inner adjustment
by a commutator proword therefore does not change the determinant.

The peripheral action used in the block substitutions is based at
infinity. For \(z=(xy)^{-1}\), put \(a=f(x,z)\),
\(b=f(y,z)\), and \(c=f(x,y)\). The two- and three-cycle
relations give \(b=ac^{-1}\). Consequently this action is
\(\operatorname{Inn}(a)\phi_c\), and its determinant is
\(B_f(X,Y)\).

The action \(F^{1,2}\) adds a fixed third generator, so its
determinant is \(B_f(X,Y)\). Similarly \(F^{2,3}\) has
determinant \(B_f(Y,Z)\). For \(F^{12,3}\), use the free
Nielsen basis \(x,a=xy,z\). The images of \(a,z\) depend only
on \(a,z\). The image of \(x\) is its conjugate by a word in
\(a,z\) with trivial abelianization. Its first Jacobian diagonal
entry is therefore one, and the remaining lower block is the
two-generator Jacobian in variables \(XY,Z\). Its determinant
is \(B_f(XY,Z)\). The change-of-basis factors cancel because
the automorphism is IA, so changing back does not alter this
determinant. The same argument gives \(B_f(X,YZ)\) for
\(F^{1,23}\).

The two pure braids \(x_{12},x_{23}\) act trivially on the
abelianization of the rank-three fiber. The determinant of their
Fox representation is consequently a continuous homomorphism from
the free base group \(\widehat F(x_{12},x_{23})\) to a
commutative unit group. It kills the commutator proword \(f\),
so the pure-braid side has determinant one. This conclusion also
holds if the inner action is inverted.

The determinant of the fourfold block product is
\[
 \frac{B_f(X,Y)B_f(XY,Z)}{B_f(X,YZ)B_f(Y,Z)}.
\]
Equating it to the determinant-one pure-braid side proves
\eqref{eq:twist-determinant-beta}. All calculations commute with
finite coefficient reduction and inverse limits, so they apply to
the full continuous actions.
\end{proof}

The four factors are strand-substitution actions on the rank-three
fiber. This description does not give a commutator factorization in
GT. The proposition identifies only the information captured by the
scalar determinant; it makes no assertion about the remaining
matrix entries.
